\chapter{Αναπαραγωγή των Αποτελεσμάτων}
\label{app:repro}

Το παράρτημα αυτό περιγράφει την πλήρη διαδικασία αναπαραγωγής των ποσοτικών αποτελεσμάτων της εργασίας: την προετοιμασία του περιβάλλοντος, την παραγωγή των πινάκων, την εκτέλεση των μετρήσεων και τους ελέγχους συνέπειας. Κάθε παραγόμενο τεκμήριο (\eng{artifact}) φέρει μεταδεδομένα προέλευσης που καταγράφουν την κατάσταση του πηγαίου κώδικα από την οποία παράχθηκε. Τα τεκμήρια αναφοράς της εργασίας προέρχονται από καθαρές δεσμευμένες καταστάσεις του κώδικα: το δηλωτικό (\eng{manifest}) των πινάκων συντεταγμένων καταγράφει τη δέσμευση (\eng{commit}) \code{dafe9ca} με καθαρό δέντρο εργασίας (\code{dirty=false}) και ένδειξη αναπαραγώγιμου σημείου ελέγχου. Τα μεταδεδομένα του γωνιακού πίνακα προτύπων (\eng{pattern database}, PDB) καταγράφουν αντίστοιχα τη δέσμευση \code{f8d9a12}. Έτσι κάθε αριθμός του κειμένου μπορεί να συνδεθεί με συγκεκριμένο πηγαίο κώδικα. Η μόνη ειδική περίπτωση είναι ο μεγάλος πίνακας αποκοπής (\eng{pruning table}) \code{h48h7}, ο οποίος διατηρείται ως υφιστάμενο δυαδικό τεκμήριο και επανατεκμηριώνεται με ρητή διαδικασία υιοθέτησης μεταδεδομένων, όπως περιγράφεται στην Ενότητα~\ref{sec:repro-h48h7}.

\section{Περιβάλλον εκτέλεσης}
\label{sec:repro-env}

Η αναπαραγωγή απαιτεί Python έκδοσης 3.10 ή νεότερης, μεταγλωττιστή C/C++ για τις εγγενείς (\eng{native}) μηχανές και πλήρη εργαλειοθήκη TeX (XeLaTeX, \code{latexmk}, \code{bibtex}) για την κατασκευή του κειμένου. Η βασική εγκατάσταση γίνεται σε εικονικό περιβάλλον:

\begin{verbatim}
python -m venv .venv
source .venv/bin/activate
python -m pip install --upgrade pip
python -m pip install -e ".[dev]"
\end{verbatim}

Το μηχάνημα αναφοράς των μετρήσεων είναι συμβατικός υπολογιστής με 8 πυρήνες και 16~GiB μνήμης. Η πλήρης δοκιμαστική σουίτα (\code{python -m pytest -q}) ολοκληρώνεται χωρίς αποτυχίες στη βασική γραμμή \code{dafe9ca}, την ίδια που καταγράφεται στο δηλωτικό των πινάκων συντεταγμένων.

\section{Σειρά αναπαραγωγής}
\label{sec:repro-order}

Η προτεινόμενη σειρά εκτέλεσης ξεκινά από τις δοκιμές, συνεχίζει με την παραγωγή των πινάκων συντεταγμένων και των πινάκων προτύπων, έπειτα με τις μετρήσεις, τον έλεγχο των αποτελεσμάτων, την αναπαραγωγή των πινάκων LaTeX και τέλος με την κατασκευή του PDF. Η σειρά αυτή εξασφαλίζει ότι το κείμενο χτίζεται πάνω σε ενημερωμένα δεδομένα: αν αποτύχουν οι δοκιμές, δεν έχει νόημα η παραγωγή νέων πινάκων, και αν αποτύχει ο επαληθευτής, τα αποτελέσματα δεν ενσωματώνονται στο κείμενο.

\begin{verbatim}
python -m pytest -q
python scripts/generate_tables.py \
  --profile thesis --seed 2026
python scripts/generate_corner_pdb.py \
  --profile thesis --seed 2026
python scripts/generate_edge_pdb.py \
  --profile thesis --seed 2026
python scripts/generate_h48_tables.py \
  --profile thesis --seed 2026 --solver h48h0
python scripts/run_3x3_optimal.py \
  --profile thesis --seed 2026
python scripts/run_benchmarks.py \
  --profile thesis --seed 2026
python scripts/generate_pocket_cube.py \
  --profile thesis --seed 2026
python scripts/verify_results.py
python scripts/generate_figures.py \
  --profile thesis --seed 2026
latexmk -xelatex -interaction=nonstopmode \
  -halt-on-error thesis/main.tex
\end{verbatim}

Όλα τα τεκμήρια αναφοράς της εργασίας προέρχονται από το προφίλ \code{thesis} με σπόρο 2026, και τα ονόματα των αρχείων τους περιέχουν και τα δύο στοιχεία. Το προφίλ \code{quick} ολοκληρώνεται γρηγορότερα και προορίζεται για δοκιμαστικούς ελέγχους κατά την ανάπτυξη· δεν παράγει τεκμήρια αναφοράς. Το προφίλ \code{stress} αποτελεί προαιρετική διαδρομή για βαθύτερη διερεύνηση, όπως οι περιπτώσεις πίεσης βάθους 25 του Κεφαλαίου~\ref{chap:experiments}.

\section{Χαρτογράφηση τεκμηρίων}
\label{sec:repro-artifacts}

Οι πίνακες συντεταγμένων γράφονται στον κατάλογο \path{data/generated/} και συνοδεύονται από το δηλωτικό \path{data/generated/table_manifest_thesis_seed_2026.json}, το οποίο καταγράφει την προέλευση κάθε πίνακα. Τα βασικά τεκμήρια αναφοράς είναι τα ακόλουθα:

\begin{verbatim}
results/processed/table_metadata_thesis_seed_2026.json
data/generated/thesis_seed_2026_corner_state_pdb.bin
results/processed/corner_pdb_metadata_seed_2026_thesis.json
results/processed/edge_pdb_metadata_seed_2026_thesis.json
data/generated/h48/thesis_seed_2026/h48h7.bin
results/processed/h48_metadata_seed_2026_thesis_h48h0.json
results/processed/h48_metadata_seed_2026_thesis_h48h7.json
results/processed/optimal_3x3_seed_2026_thesis.json
results/raw/benchmarks_seed_2026_thesis.jsonl
results/processed/summary_seed_2026_thesis.json
results/processed/benchmarks_seed_2026_thesis.csv
results/raw/pocket_cube_distribution_seed_2026_thesis.json
results/processed/pocket_cube_summary_seed_2026_thesis.json
\end{verbatim}

Οι ωμές γραμμές μετρήσεων και η πλήρης κατανομή του κύβου 2×2×2 βρίσκονται στον κατάλογο \path{results/raw/}, ενώ οι επεξεργασμένες συνόψεις και τα μεταδεδομένα στον κατάλογο \path{results/processed/}. Το σχήμα των γραμμών μετρήσεων τεκμηριώνεται στο Παράρτημα~\ref{app:schema}. Οι πίνακες LaTeX και τα αρχεία δεδομένων των σχημάτων παράγονται από τα επεξεργασμένα αποτελέσματα με το \code{scripts/generate_figures.py} και εισάγονται στο κείμενο ως παραγόμενα αρχεία, χωρίς χειροκίνητη αντιγραφή αριθμών.

\section{Αναμενόμενοι χρόνοι και απαιτήσεις μνήμης}
\label{sec:repro-cost}

Οι ακριβείς χρόνοι εκτέλεσης διαφέρουν ανά μηχάνημα και φορτίο· οι ετικέτες κατάστασης των αποτελεσμάτων οφείλουν όμως να παραμένουν ίδιες για το ίδιο προφίλ, εφόσον δεν αλλάξουν ο κώδικας ή τα όρια. Ενδεικτικά, στο μηχάνημα αναφοράς η παραγωγή του γωνιακού πίνακα προτύπων διήρκεσε 9.88015 δευτερόλεπτα, ενώ η παραγωγή του πίνακα \code{h48h7} διήρκεσε 2359.7427 δευτερόλεπτα με 8 νήματα. Ο πίνακας \code{h48h7} καταλαμβάνει 3\,793\,842\,344 bytes στον δίσκο, και η χρήση του σε μηχάνημα με 16~GiB μνήμης είναι εφικτή αλλά αισθητά εξαρτημένη από το διαθέσιμο περιθώριο. Η πρώτη ψυχρή φόρτωσή του είναι ακριβή, ενώ οι πίνακες μεγαλύτερων επιπέδων (\code{h48h8} και άνω) υπερβαίνουν τις δυνατότητες του μηχανήματος αναφοράς. Η εγγενής εξαντλητική απόδειξη βελτιστότητας του superflip, η ακριβότερη μεμονωμένη εκτέλεση της εργασίας, απαιτεί περίπου 87 λεπτά στο ίδιο μηχάνημα (Κεφάλαιο~\ref{chap:experiments}).

\section{Έλεγχοι συνέπειας}
\label{sec:repro-checks}

Το σενάριο \code{scripts/verify_results.py} εκτελείται μετά από κάθε παραγωγή μετρήσεων: επιβεβαιώνει το σχήμα των αποτελεσμάτων, επαναλαμβάνει την επαλήθευση κάθε δηλωμένης λύσης και διασταυρώνει τις ρηχές ακριβείς αποστάσεις με εξαντλητική BFS (Παράρτημα~\ref{app:schema}). Το σενάριο \code{scripts/thesis_audit.py} ελέγχει συμπληρωματικά τη χρονική συνέπεια των τεκμηρίων: εντοπίζει απαιτούμενα παραγόμενα αρχεία που λείπουν, είναι κενά ή είναι παρωχημένα (\eng{stale}), δηλαδή παλαιότερα από την πηγή τους. Ο έλεγχος αυτός δεν αποδεικνύει από μόνος του επιστημονική ορθότητα, εντοπίζει όμως τα συνηθέστερα λάθη αναπαραγωγής: για παράδειγμα, αν αλλάξει μια σύνοψη αποτελεσμάτων χωρίς να αναπαραχθεί ο αντίστοιχος πίνακας LaTeX, ή αν λείπει ο γωνιακός πίνακας προτύπων, κάποιος πίνακας ακμών ή ο πίνακας του κύβου 2×2×2, ο έλεγχος αποτυγχάνει και υποδεικνύει το ελλιπές τεκμήριο.

\section{Η εξαίρεση του πίνακα h48h7}
\label{sec:repro-h48h7}

Ο πίνακας αποκοπής \code{h48h7} αποτελεί τη μοναδική τεκμηριωμένη εξαίρεση στη byte προς byte αναπαραγωγή των δυαδικών πινάκων. Η πολυνηματική παραγωγή του ήταν, κατά την περίοδο παραγωγής, μη ντετερμινιστική, με συνέπεια η εντολή \code{python scripts/generate_h48_tables.py --oracle} να παράγει λειτουργικά ισοδύναμο πίνακα, του οποίου όμως το άθροισμα ελέγχου μπορεί να διαφέρει από εκείνο του διατηρημένου τεκμηρίου. Για τον λόγο αυτόν, η τελική διαδικασία δεν αλλάζει σιωπηρά το άθροισμα ελέγχου ούτε ισχυρίζεται νέα byte προς byte παραγωγή. Αντίθετα, η εντολή \code{python scripts/generate_h48_tables.py --oracle --adopt-existing-table-metadata} εκτελείται από καθαρή δεσμευμένη κατάσταση του κώδικα, επαληθεύει τον υπάρχοντα πίνακα με εγγενή έλεγχο αποδοχής, υπολογίζει εκ νέου το SHA-256, καταγράφει καθαρή προέλευση για την πράξη υιοθέτησης και διατηρεί την προηγούμενη προέλευση στα πεδία \code{adoption_previous_*}. Στη συνέχεια αναπαράγονται τα τεκμήρια πιστοποίησης που ενσωματώνουν τα μεταδεδομένα του πίνακα. Όλες οι βαθιές ακριβείς γραμμές που στηρίζονταν αποκλειστικά στον \code{h48h7} έχουν επιβεβαιωθεί ανεξάρτητα από το εξωτερικό Nissy 2.0.8. Η πλήρης ανάλυση του ευρήματος, της οριοθέτησής του και της απόφασης διατήρησης δίνεται στο Κεφάλαιο~\ref{chap:validation}.

\section{Πλήρης καταγραφή των ελέγχων συμβολαίου του μαντείου H48}
\label{sec:repro-contract}

Για λόγους πληρότητας, ο Πίνακας~\ref{tab:h48-contract-full} παραθέτει την πλήρη, μη επιμελημένη καταγραφή των ελέγχων συμβολαίου (\eng{contract checks}) της διαδρομής μαντείου (\eng{oracle}) H48, όπως παράγεται από την εντολή:
\begin{verbatim}
python scripts/generate_h48_oracle_contract.py \
  --profile thesis --seed 2026 --solver h48h7
\end{verbatim}
Η ερμηνεία των κρισιμότερων ελέγχων, ιδίως της διάκρισης ανάμεσα στο εμπειρικό σώμα μετρήσεων και στο πεδίο που δηλώνει ότι δεν υπάρχει απόδειξη ταχύτητας για κάθε δυνατή κατάσταση, δίνεται στο Κεφάλαιο~\ref{chap:experiments}. Οι τιμές του πίνακα προέρχονται αυτούσιες από το αποθηκευμένο τεκμήριο και περιλαμβάνουν εσωτερικά αναγνωριστικά ελέγχων στην αγγλική. Σημειώνεται ότι η γραμμή \code{Resident speedup} αποτυπώνει την παλαιά λογιστική χωρίς κοινή προθέρμανση· η δίκαιη μέτρηση της αντίστοιχης επιτάχυνσης δίνεται στο Κεφάλαιο~\ref{chap:experiments}.

\begin{table}[htbp]\centering
{\small
\begin{tabular}{@{}p{0.54\linewidth}p{0.34\linewidth}@{}}
\hline
Check & Value \\
\hline
Fast optimal oracle API & True \\
All-state exact contract & True \\
Empirical fast corpus & True \\
Every-state fast runtime proof & False \\
Fast API max seconds & 2.664585 \\
Resident hard max seconds & 77.822389 \\
Resident speedup & 40.16 \\
Nissy optimal stress max seconds & 23.734412 \\
Nissy-core direct max seconds & 2.418875 \\
Portfolio Nissy-first max seconds & 10.097333 \\
Portfolio state recovery max seconds & 11.431923 \\
Portfolio direct state max seconds & 2.399119 \\
Race oracle max seconds & 2.668843 \\
Race direct state seconds & 3.993046 \\
Resident race max seconds & 5.829176 \\
Resident race direct seconds & 4.46796 \\
Universal direct state seconds & 5.189862 \\
Universal H48 symmetry seconds & 3.681836 \\
Universal H48 batch seconds & 18.413704 \\
Universal CLI max seconds & 3.712856 \\
Universal CLI broader seconds & 48.364612 \\
Universal CLI adaptive seconds & 10.477171 \\
Universal CLI expanded seconds & 22.674068 \\
Universal CLI live no-shortcuts seconds & 4.615232 \\
Universal CLI live no-shortcuts broader seconds & 15.582677 \\
Universal CLI known-distance 19 seconds & 227.559559 \\
Universal CLI known-distance 20 seconds & 508.871251 \\
Portfolio superflip fallback seconds & 135.679618 \\
Portfolio superflip cache seconds & 0.012885 \\
Certificate symmetry cases & 736 \\
Expanded certificate symmetry cases & 1518 \\
Learned certificate replay seconds & 0.00789 \\
Live symmetry batch max seconds & 1.651504 \\
Nissy optimal table bytes & 2465897307 \\
h48h8 probe status & timed\_\allowbreak{}out \\
h48h8 probe allocated bytes & 6026702848 \\
h48h8 detached status & detached\_\allowbreak{}process\_\allowbreak{}not\_\allowbreak{}alive\_\allowbreak{}no\_\allowbreak{}trusted\_\allowbreak{}table \\
h48h8 detached waits & 160 \\
\hline
\end{tabular}
}

\caption{Πλήρης καταγραφή των ελέγχων συμβολαίου του μαντείου H48 (\code{h48h7}), όπως παράγεται από το \code{scripts/generate_h48_oracle_contract.py}.}
\label{tab:h48-contract-full}
\end{table}

\section{Περιορισμοί αναπαραγωγής}
\label{sec:repro-limits}

Η αναπαραγωγή εξαρτάται από διαθέσιμο περιβάλλον Python και εργαλειοθήκη TeX. Οι μετρήσεις χρόνου εκτέλεσης είναι ενδεικτικές του μηχανήματος αναφοράς και δεν αναμένεται να συμπέσουν αριθμητικά σε άλλο υλικό. Η διαθεσιμότητα του εξωτερικού πακέτου \code{kociemba} \cite{muodovKociembaPkg} είναι προαιρετική εξάρτηση: αν το πακέτο απουσιάζει, η διαδρομή του προσαρμογέα δηλώνεται μη διαθέσιμη χωρίς να επηρεάζονται οι υπόλοιπες διαδρομές. Για τον λόγο αυτόν, τα κύρια εσωτερικά τεκμήρια της εργασίας είναι οι εγγενείς διαδρομές και ο ανεξάρτητος επαληθευτής, των οποίων η συμπεριφορά δεν εξαρτάται από εξωτερικές εγκαταστάσεις.
